\documentclass{report}
\usepackage{amsmath,amssymb}
\setlength{\parindent}{0mm}
\setlength{\parskip}{1em}
\begin{document}
\begin{center}
\rule{6in}{1pt} \
{\large J. David Moulton \\
{\bf Multiphysics Frameworks and Scalable Solvers for Environmental Applications}}

Applied Mathematics and Plasma Physics \\ Los Alamos National Laboratory \\ Los Alamos \\ NM 87545
\\
{\tt moulton@lanl.gov}\\
Ethan Coon\\
Gianmarco Manzini\\
Markus Berndt\\
Rao Garimella\\
Konstantin Lipnikov\\
Scott Painter\end{center}

Modeling and simulation are playing an increasingly critical role in
understanding and predicting climate impacts and feedbacks in terrestrial
systems. Managing the complexity of these process-rich integrated
hydrologic and biogeochemical models requires flexible software designs
that enable exploration of model features and model coupling. In
addition, flexibility in meshing and robust discretization techniques are
required to capture topographic features, such as hill slopes and rivers,
and suburface stratigraphy. This combination of process-rich applications
with advanced discretizations creates significant challenges for
efficient solvers, and efficient time-evolution.

In this talk we highlight a flexible and extensible approach to
multiphysics frameworks for these applications that specifies interfaces
for coupled processes, automates weak coupling and supports strong
coupling strategies to manage this complexity. We use a prototype
implementation this multiphysics framework, dubbed Arcos, to support both
model and algorithm development for environmental applications in the
open-source codes Amanzi and the Arctic Terrestrial Simulator (ATS).
Amanzi provides infrastructure for general unstructured polyhedral
meshing with a flexible operator-based implementation of the Mimetic
Finite Difference method as well as two-point flux Finite Volume method.
Amanzi is used for subsurface flow and reactive transport, while ATS
includes capabilities for coupled non-isothermal surface/subsurface flow.

Here we examine the increasingly important problem of coupled
surface/subsurface flow where we use a diffusive wave approximation for
surface flow, and a Richards equation for subsurface flow as a testbed
for discretizations and solvers. In both equations, the coefficients of
the div/grad operator are strongly nonlinear, and the surface and
subsurface flows are strongly coupled. The challenges of designing a
preconditioner for Krylov-based nonlinear solvers with first-order
backward Euler time discretization are discussed. For example, upwinding
of the nonlinear coefficient is required to evolve the system
efficiently, but breaks the symmetry of the operator on general meshes.
Furthermore, as the ponded height of the surface water goes to zero
(e.g., water runs off or infiltrates into the subsurface) the nonlinear
coefficient goes to zero, making standard formulations that involve its
inverse invalid. Finally, the system is strongly coupled through
continuity of both pressure and flux at the surface/subsurface interface,
and adaptive time-stepping is required to evolve the system efficiently.
We discuss the performance of the nonlinear and multilevel linear solvers
(e.g., Hypre BOOMERAMG), on this fully coupled system at various times in
the simulation, and the hueristics used to control the dynamic time-step.
We show results for several benchmark problems, as well as physically
relevant, large-scale simulation of rainfall on arctic tundra based upon
LIDAR data from Barrow, Alaska.


\end{document}
